\section{运动学}
\lsn{1}运动学的五个参数：\lunl{3em}、\lunl{3em}、\lunl{3em}、\lunl{3em}、\lunl{3em}。

\lsn{2}根据\(v - t\)图像得出匀变速直线运动的公式：

\begin{tikzpicture}[>=stealth,font=\footnotesize]
    \draw[|<->|,thin,gray] (3.25,2.5)--node[fill=white]{ \(\Delta v\)}(3.25,1);
    \filldraw[pattern=north east lines,draw=white,pattern color=red!30] (0,1)--(3,2.5)--(3,1)--(0,1);
    \filldraw[pattern=north west lines,draw=white,pattern color=red!30] (0,0) rectangle(3,1);
    \draw[->,blue] (0,0)--(5,0)node[below]{ \(t\)};
    \draw[->,blue] (0,0)node[below left]{\(O\)}--(0,4)node[right]{ \(v\)};
    \draw[magenta] (0,1)--(4,3)node[below right,black]{ \(a\)};
    \draw[dashed,cyan] (0,2.5)--(3,2.5)--(3,0) (0,1)--(3,1) (0,1.75)--(1.5,1.75)--(1.5,0);
    \node[blue] at(0,1)[left]{ \(v_0\)};
    \node[blue] at(0,2.5)[left]{ \(v_t\)};
    \node[blue] at(3,0)[below]{ \(t\)};
    \node[blue] at(0,1.75)[left]{ \(\overline{v}\)};
    \node[blue] at(1.5,0)[below]{\(\sfrac{t}{2}\)};
    \draw[|<->|,gray,thin] (0,-.8)--node[fill=white]{ \(\Delta t\)}(3,-.8);
  \end{tikzpicture}


\lsn{3}证明：匀变速直线运动中，\(\overline{v} = \frac{v_0 + v}{2} = v_{\frac{t}{2}}\)

\vspace*{3cm}
\lsn{4}运动的物体什么情况可以看成质点？什么情况不可以看成质点？

\vspace*{3cm}
\lsn{5}简述电磁打点计时器与电火花打点计时器的异同

\vspace*{3cm}
\lsn{6}推导：自设情景并导出\(\Delta x = aT^2\)

\newpage
\lsn{7}写出初速度为0的匀加速直线运动的六个比例关系

\vspace*{5cm}
\lsn{8}写出自由落体运动的定义及公式

\vspace*{4cm}

\lsn{9}写出竖直上抛运动的公式

\vspace*{4cm}
\lsn{10}怎样理解竖直上抛运动的对称性？

\vspace*{3cm}
\lsn{11}在表中填出各个图像中不同元素的物理含义，没有物理意义的画横线

\begin{table}[htbp]
  \begin{tabular}{l|c|c|c}
    \hline
    &\(x -t\)图像&\(v - t\)图像&\(a - t\)图像\\
    \hline
    斜率&&&\\
    \hline
    面积&&&\\
    \hline
    交点&&&\\
    \hline
    零点&&&\\
    \hline
    拐点&&&\\
    \hline
  \end{tabular}
\end{table}

\newpage
\lsn{12}定出下列图像中初速度\(v_0\)、加速度\(a\)，并给出过程

\begin{tikzpicture}[>=stealth,font=\footnotesize]
  \draw[->,blue] (0,0)node[below left]{\(O\)}--(4,0)node[below]{\(t\)};
  \draw[->,blue] (0,0)--(0,3)node[right]{\(\overline{v}\)};
  \draw[magenta] (0,.8)node[left,blue]{\(b\)}--(3,2.5)node[below,blue]{\(k\)};
  \begin{scope}[xshift=5cm]
    \draw[->,blue] (0,0)node[below left]{\(O\)}--(4,0)node[below]{\(\frac{1}{t}\)};
  \draw[->,blue] (0,0)--(0,3)node[right]{\(\frac{x}{t^2}\)};
  \draw[magenta] (0,.8)node[left,blue]{\(b\)}--(3,2.5)node[below,blue]{\(k\)};
  \end{scope}
  \begin{scope}[xshift=10cm]
    \draw[->,blue] (0,0)node[below left]{\(O\)}--(4,0)node[below]{\(x\)};
  \draw[->,blue] (0,0)--(0,3)node[right]{\(v^2\)};
  \draw[magenta] (0,.8)node[left,blue]{\(b\)}--(3,2.5)node[below,blue]{\(k\)};
  \end{scope}
\end{tikzpicture}

\vspace*{3cm}
\lsn{13}下图中每两点间时间均为\(T\)，则

\begin{tikzpicture}
		\coordinate (o) at(0,0);
		\draw ($(o)+(-.3,.5)$)--++(10.2,0) ($(o)+(-.3,-.5)$)--++(10.2,0)coordinate(a) ($(o)+(-.3,.5)$)--++(-.1,-.1)--++(.2,-.7)--++(-.1,-.2) (a)--++(.1,.1)--++(-.2,.7)--++(.1,.2);
		\node at($(a)+(-.7,1.2)$){\scriptsize 单位：\,\qty{}{cm}};
		\foreach \x/\n in {0/0,.6/1,1.6/2,3/3,4.8/4,7/5,9.6/6}{
			\coordinate (\n) at(\x,0);
			\fill (\n)node[above]{\scriptsize \(\n\)} circle(.03);}
		\foreach \y/\x/\n in {-.7/.6/L_1,-1.1/1.6/L_2,-1.5/3/L_3,-1.9/4.8/L_4,-2.3/7/L_5,-2.7/9.6/L_6}
			\draw[stealth-stealth,thin] (0,\y)--node[below]{\scriptsize \(\n\)}(\x,\y);
		\foreach \x/\y in {0/-2.8,.6/-.8,1.6/-1.2,3/-1.6,4.8/-2.0,7/-2.4,9.6/-2.8}
			\draw[gray] (\x,-.1)--(\x,\y);

      \foreach \x/\y/\n in {0/.6/x_1,.6/1.6/x_2,1.6/3/x_3,3/4.8/x_4,4.8/7/x_5,7/9.6/x_6}
			\draw[stealth-stealth,thin] (\x,0)--node[fill=white]{\scriptsize \(\n\)}(\y,0);
	\end{tikzpicture}%%始.6,差.4

\(v_1 =\)\lunl{3cm}\qquad \(v_4 =\)\lunl{3cm}

\(v_0 =\)\lunl{3cm}\qquad \(v_6 =\)\lunl{3cm}

\(a =\)\lunl{6cm}（四个数据）

\(a =\)\lunl{6cm}（五个数据）

\(a =\)\lunl{6cm}（六个数据）
